
\documentclass[11pt]{article}
\usepackage[utf8]{inputenc}
\usepackage[spanish]{babel}
\usepackage{amsmath}
\usepackage{geometry}
\geometry{margin=1in}

\title{Documentación: Clase 1 Series de Tiempo}
\date{}

\begin{document}
\maketitle

\section{Objetivo}
Introducir los conceptos fundamentales del análisis de series de tiempo.

\section{Temas Principales}
\begin{itemize}
\item Definición de serie de tiempo.
\item Componentes: tendencia, estacionalidad y ruido.
\item Estacionariedad.
\item Autocorrelación.
\item Visualización y análisis exploratorio.
\end{itemize}

\section{Importancia}
Las series de tiempo permiten analizar datos observados secuencialmente y realizar pronósticos futuros.

\section{Conclusión}
Esta práctica introduce las bases necesarias para comprender modelos más avanzados como ARIMA, SARIMA y GARCH.

\end{document}
